% ch00-howto — reading guide (unnumbered front-matter chapter)

This document is the master reference for the \code{hawkes-sim} modeling
effort: every future code modification in \modrepo{} should be traceable to a
numbered decision, assumption, or protocol stated here. It supersedes the
earlier working notes (\code{THEORY.md}, \code{MATH\_NOTES.md},
\code{STRATEGY.md}) as the single source of truth, while carrying forward
their results wherever those results survived scrutiny; the numerical
verdicts inherited from the synthetic rehearsal campaigns are collected in
\cref{app:results} so that the main text can cite them without re-deriving
them.

The document is organized around one organizing question --- \emph{what is
valuable to know?} --- posed in \cref{ch:question} and answered differently
by every subsequent part. \Cref{ch:data} fixes, once, what the
PitchBook/FactSet data actually delivers and the six ways observation
corrupts it. \Cref{ch:foundations,ch:poisson,ch:hawkes,ch:censoring,ch:gof}
build the continuous-time machinery: point-process foundations, Poisson and
exponential-Hawkes models with covariate-driven excitation and inhibition,
the treatment of untrusted and censored deal times, and the goodness-of-fit
discipline. \Cref{ch:stale} confronts the covariate staleness that PitchBook
imposes. \Cref{ch:buckets,ch:ranking} develop the discrete-time bucket
formulation solved with gradient boosting and the ranking view evaluated by
recall at $K$. \Cref{ch:universe} handles deals arriving from companies
outside the tracked universe. \Cref{ch:decision,ch:roadmap} close the loop:
from model outputs to investment decisions, and from decisions to code.

Three reading conventions. First, \emph{assumptions are contracts}: every
model result is stated under numbered assumptions, and the text says
explicitly which assumptions are believed, which are convenient fictions,
and which are known to be false in the data (with the consequence of the
violation). Second, \emph{Design decision} boxes record commitments --- a
choice among alternatives, with the rationale and the trigger that would
reverse it. Third, \emph{Code hooks} boxes at the end of each chapter map
the mathematics onto modules of \modrepo{}, existing or to be written; these
boxes, together with \cref{ch:roadmap}, are the implementation contract.

Proofs are given in the running text when they are short and instructive;
longer arguments are deferred to the ends of chapters. Citations are
restricted to work we have actually relied on --- either for a theorem, a
method we implement, or a warning we heed --- and the bibliography is
annotated in that spirit throughout the text rather than in a separate
rubric.
